\section{Discussion}
\label{sec:discussion}

\subsection{Comparison Table}
\label{sec:compare-table}

\Cref{tab:scheme-compare} situates \TW{} against representative authenticated
encryption schemes along the axes that distinguish it: the underlying primitive,
whether the mode is parallelizable, the key, nonce, and tag sizes, associated data
support, nonce-misuse resistance, and the published committing result or attack
most relevant to the scheme. The rows are representative rather than
exhaustive: sponge/duplex relatives, a Farfalle-based parallel design, AES
baselines, a misuse-resistant AES baseline, and a generic committing transform.
The comparison is structural, not a total security ranking; \TW{}'s own
concrete bounds and the regimes in which each term dominates are treated in
\Cref{sec:bound-discussion}, and its measured throughput against
hardware-accelerated AES-128-GCM in \Cref{sec:compare-perf}. The committing
column should be read within the named notions only, and a negative entry there
does not assert a confidentiality or integrity break of the underlying AEAD. It
records \TW{}'s $128$-bit \textsf{CMTDs-4} bound derived from tag hash
security, Ascon's roughly $64$-bit committing bound \cite{KSW23}, published
low-cost attacks where known \cite{KSW23,CR22,ADGKLS22}, and the
hash-collision level inherited by \textsf{CTX}.
The table does not separately track compact (tag-only) commitment, since
that property has not been systematically analyzed for the comparison
schemes.

\begin{table}[t]
\centering
\footnotesize
\caption{Structural comparison of \TW{} with representative AEADs. Sizes are in
bytes; AES-GCM-SIV also admits a $32$-byte key, and Kravatte-SAE's key and
nonce lengths are parameters (keys of up to $40$ bytes, arbitrary-length
nonces). ``Misuse'' is nonce-misuse
resistance in the sense of \cite{RS06}; ``Commit'' records a published
committing result or attack in the relevant committing notion, with ``broken''
denoting a published low-cost committing attack, ``hash CR'' denoting the
collision-resistance level of the hash used by the transform, and ``open''
marking a scheme with no published committing analysis. A dash means the entry
depends on the base scheme. These entries are notion-specific summaries, not
broad AE security claims. \TW{}'s concrete bounds appear in
\Cref{sec:bound-discussion}.}
\label{tab:scheme-compare}
\begin{tabular}{@{}llccccc@{}}
\toprule
Scheme & Primitive & Par. & K/U/T & AD & Misuse & Commit \\
\midrule
\TW{}                        & \Keccakp   & yes & $32/32/32$ & yes & no  & CMTDs-4 \\
Ascon-128a \cite{DEMS21}     & Ascon-$p$ ($320$ b) & no  & $16/16/16$ & yes & no  & $\approx64$ b \\
Xoodyak \cite{DHPVV20}       & Xoodoo ($384$ b)   & no  & $16/16/16$ & yes & no  & $2^{17}$ atk. \\
Kravatte-SAE \cite{BDHPVV17} & \Keccak[1600, 6]   & yes & $\leq 40$/var./$16$ & yes & no & open \\
AES-GCM \cite{Dwo07}         & AES                & yes & $16/12/16$ & yes & no  & broken \\
AES-GCM-SIV \cite{GLL17}     & AES                & yes & $16/12/16$ & yes & yes & broken \\
nAE $+$ \textsf{CTX} \cite{CR22} & any nAE        & ---  & ---        & yes & --- & hash CR \\
\bottomrule
\end{tabular}
\end{table}

\subsection{Comparison with Prior Sponge AEADs}
\label{sec:compare-sponge}

\TW{} belongs to the single-permutation duplex AEAD family introduced in
\cite{BDPV11}. Representative lightweight members of that family, Ascon
\cite{DEMS21} and Xoodyak \cite{DHPVV20}, target $128$-bit security on small
permutations ($320$ and $384$ bits) with a strictly serial duplex: each call
depends on the previous one, so there is no mode-level parallelism to exploit.
\TW{} instead pays for a larger permutation, \Keccakp, and a $256$-bit
capacity---the same round count as KangarooTwelve \cite{BDPVVV16}---in exchange
for a tree that processes leaf chunks in parallel and for the headroom that
meets the $2^{-128}$ target at the stated work cap. The cost of this choice is
visible in \Cref{sec:performance}: on short, single-leaf inputs \TW{} is slower
per byte than a compact serial duplex would be, and its advantage appears only
once a message spans enough leaves to fill the vector units.

\TW{} reuses KangarooTwelve's final-node/leaf split but not its Sakura coding
\cite{BDPV14}. KangarooTwelve uses Sakura to make the final node parseable as a
tree-hash transcript \cite{BDPVVV16}. In \TW{}, the analogous chaining values
are typed leaf tags: each leaf duplex is initialized with its leaf prefix and
index, and the root absorbs leaf tags in index order followed by the leaf count
(\Cref{sec:tw128,sec:domain-sep}). Sakura's structural role is therefore
handled by mode initialization and domain separation rather than by a tree
coding. The transcript separation lemmas make this explicit
(\Cref{lem:transcript-inj,cor:output-sep}); the hash collision bound sees the
leaf encoding only through the explicit leaf tag collision term of
\Cref{thm:coll}.

Keyak \cite{BDPVV16}, whose Motorist mode is already parallelizable, is the
closest prior duplex AEAD in spirit; \Cref{sec:related} contrasts its
wire-visible, encryptor-and-decryptor-shared degree of parallelism with \TW{}'s
run-time leaf choice.
Strobe \cite{Ham17} is a serial sponge framework aimed at whole protocols rather
than at a standalone AEAD, and like Ascon and Xoodyak it does not parallelize.

Farfalle \cite{BDHPVV17} leaves the duplex setting altogether: it is a
parallel pseudorandom function whose compression layer sums
masked-and-permuted input blocks into an accumulator and whose expansion
layer generates output from a rolling state, with authenticated encryption
supplied by the Farfalle-SAE and Farfalle-SIV modes on top
(\Cref{sec:related}). Its Kravatte instance runs \Keccak[1600, 6], and the
parallelism available to an implementation is arbitrary, as in \TW{}. The two
designs price their permutations for different adversary models. Farfalle's
six-round compression layer is analyzed in a PRF setting where inputs are
randomized by secret masks and an adversary does not see both the input and
output of the same permutation call \cite{BDHPVV17}. \TW{}'s tag-hash claims
are adversarial-initialization claims: in the hash and committing games, the
adversary may choose the key as part of the input being hashed. \TW{} therefore
keeps the strict duplex object, uses the hash-grade round count discussed in
\Cref{sec:permutation-cryptanalysis}, and takes its parallelism from the tree
topology rather than from a parallel PRF compression layer.

Committing security is not unique to \TW{} in this family: a dedicated analysis
of the sponge-based NIST lightweight finalists proves Ascon committing while
breaking others \cite{KSW23}, and some constructions, such as Strobe and the
Farfalle modes, have not been examined in that setting at all. What
distinguishes \TW{} is that the authentication tag is analyzed directly as a
hash output for the full input. Collision, second-preimage, preimage, and
everywhere preimage bounds for that tag then yield the compact commitment
statements, including the $128$-bit full-input \textsf{CMTDs-4} bound, as
mode-native consequences without leaving the single-permutation duplex setting.

\subsection{Comparison with Committing AEAD Constructions}
\label{sec:compare-commit}

The dominant route to committing AEAD is a generic transform over a
non-committing base scheme. \cite{BH22} give the \textsf{UtC}, \textsf{RtC}, and
\textsf{HtE} transforms, which add key- or context-commitment to an arbitrary
AEAD, and \cite{CR22} give \textsf{CTX}, which makes any nonce-based AE scheme
commit to its full context at the cost of a single hash over a short string and
with no ciphertext expansion. These transforms are inexpensive and broadly
applicable; \textsf{CTX} in particular is nearly free. It is nevertheless a
committing-AE transform, not a tag-hash transform: it binds the transformed
ciphertext, but does not make the tag a standalone hash output for the full
input. \TW{} differs not by undercutting their cost but by integration. The
same permutation provides confidentiality, integrity, and the tag hash output;
no separate commitment primitive, key-derivation step, or auxiliary hash is
invoked. The multi-input full-commitment bound follows from the same root tag
collision bound used for hash security (\Cref{sec:root-tag-hash-security}), and
\Cref{thm:coll} applies even when the ciphertext bodies used as openings
differ, up to the additional leaf tag birthday term. This places \TW{} alongside
the dedicated committing primitives---the compactly committing AEAD of
\cite{GLR17} and the encryptment of \cite{DGRW18}---which likewise build
commitment in rather than bolt it on, but which were designed for message
franking rather than as general nonce-based AEADs. The contrast with transformed
conventional schemes is sharpened by the attacks of \Cref{sec:related}: base
schemes such as GCM and OCB are not committing on their own \cite{CR22}, and a
transform is the only way to make them so.

Closest to \TW{} here is the recent work of \cite{DHMVV24}, which builds
nonce-based AE directly on the SHA-3 functions \textsf{SHAKE} and
\textsf{TurboSHAKE}---the latter using the same round-reduced \Keccak{}
permutation as \TW{}---and likewise achieves \textsf{CMT-4} natively, from the
collision resistance of the underlying function rather than from a transform. The
two designs make different structural choices: \cite{DHMVV24} targets
session-based communication, chaining a stream of messages through intermediate
tags and generalizing the keyed duplex into a \emph{duplex cipher}, whereas \TW{}
is a single two-level tree whose independent leaves let the implementation
choose parallelism at runtime with constant memory (\Cref{sec:related,sec:impl}).
Because a session
chain processes each message serially, \TW{}'s principal advantage is throughput
on large messages (\Cref{sec:performance}): its parallel leaves fill vector units
a serial duplex leaves idle, at the cost of the session support \cite{DHMVV24}
provides and \TW{} does not. Both demonstrate
that committing AE needs no add-on transform once it is built on a
well-scrutinized \Keccak{} permutation.

\subsection{Cryptanalysis of the Underlying Permutation}
\label{sec:permutation-cryptanalysis}

Every bound in this paper models \Keccakp{} as a uniformly random permutation
(\Cref{sec:security-model}). That modeling choice is supported by
cryptanalytic scrutiny rather than by proof, so this subsection summarizes the
third-party record and the safety margin it indicates. The Keccak round
function has been unchanged since 2008, and \TW{} uses the permutation at the
round count introduced by KangarooTwelve \cite{BDPVVV16} and later exposed
directly as TurboSHAKE \cite{BDHPVVV23}, whose designers maintain a survey of
third-party cryptanalysis. TurboSHAKE128 shares \TW{}'s rate and capacity
($r = 1344$, $c = 256$) as well as its permutation, so attacks on it and on
round-reduced SHAKE128 transfer to \TW{}'s geometry directly.

In the unkeyed setting, the deepest collision attacks on any SHA-3 or SHAKE
instance reach $6$ of \TW{}'s $12$ rounds: a SAT-aided attack on $6$-round
SHAKE128 with time $2^{123.5}$ \cite{GLST22}, and internal-differential
attacks reaching $6$ rounds of SHAKE256 and $5$ rounds of SHA3-384
\cite{ZHL24}. Preimage attacks reach $5$ rounds \cite{ZHL25}. Against the
collision- and preimage-style attacks most relevant to the root tag hash
claims of \Cref{sec:root-tag-hash-security}, the permutation therefore
retains a margin of $6$ of $12$ rounds---the assessment under which the
KangarooTwelve round count was chosen and later reaffirmed
\cite{BDPVVV16,BDHPVVV23}. The structural distinguisher reaching the most
rounds of the permutation itself, the SymSum property on $9$ rounds reported
in the survey of \cite{BDHPVVV23}, has, like the zero-sum distinguishers
before it, not been extended to an attack on any sponge or duplex mode.

The keyed setting is closer to \TW{}'s AE games: a hidden key, with
adversary-controlled nonces and messages entering through the duplex rate.
The attacks reaching the most rounds here are cube and conditional cube
attacks. Dinur et al.~break up to $9$ rounds of Keccak-based stream ciphers
and MACs faster than generic search, with keystream prediction on $8$ rounds
at time and data $2^{128}$ and on $9$ rounds (of a $512$-bit-key variant) at
$2^{256}$ \cite{DMPSS15}. With MILP-optimized conditional cubes,
nonce-respecting key recovery reaches $8$ of Lake Keyak's $12$ rounds at time
$2^{71}$ for $128$-bit keys, $9$ rounds at $2^{137}$ for $256$-bit keys, and
$9$ of KMAC256's $24$ rounds at $2^{147}$ \cite{SGLS18}. Lake Keyak is a
keyed duplex over the same $1600$-bit permutation whose initialization, like
\TW{}'s root and leaf initializations, absorbs the key and nonce ahead of
nonce-varied cube queries, so these results are the closest published analog
to \TW{}'s keyed setting. They leave a margin of $3$ to $4$ of $12$ rounds
against attacks whose complexities start at $2^{71}$.

In sum, the known attack frontier stands at $6$ of $12$ rounds for collisions
and preimages and at $8$ to $9$ rounds for keyed-mode key recovery and
structural distinguishers. \TW{} uses \Keccakp{} at the same round count and
with the same rate/capacity split as KangarooTwelve and TurboSHAKE128, so its
margin assessment tracks theirs, and any future advance against the
permutation applies to all three alike. What \TW{} adds at the mode level and
these analyses do not cover is the tree structure itself---in particular the
family of leaf initializations sharing $(K, U)$ and differing only in the
leaf index---and we consider dedicated cryptanalysis of that structure a
natural target for future work.

\subsection{Limitations and Open Problems}
\label{sec:limitations}

\Cref{sec:security-model} states the formal scope of the security claims; this
subsection discusses the consequences of the excluded settings and the
corresponding open problems.

\paragraph{Nonce misuse.}
As scoped in \Cref{sec:security-model}, \TW{}'s AE and mu-ind bounds require
nonce-respecting encryption queries, and the mode is not misuse-resistant: a
repeated nonce reuses leaf keystream and breaks confidentiality, as for any
online duplex stream cipher. Schemes such as AES-GCM-SIV \cite{GLL17}, in the
deterministic/misuse-resistant lineage of \cite{RS06}, tolerate nonce
repetition at the price of a second pass. A misuse-resistant variant of \TW{}
is left open.

\paragraph{Unverified plaintext release.}
The \textsf{INT-RUP} theorem of \Cref{thm:int-rup} covers integrity when an
adversary can obtain plaintext-computing oracle outputs before verification, in
the release-of-unverified-plaintext setting of \cite{ABLMMY14}. This is an
integrity-only statement. It does not give plaintext awareness or privacy under
release of unverified plaintext: \TW{} is online and single-pass, and exposing
unverified stream plaintext can reveal keystream information. Application
behavior on unverified plaintext is likewise outside the provable model.

\paragraph{Conservative lift.}
Per-term tightness is itemized in \Cref{sec:loss-terms}; the flat-sponge terms
are conservative because our proof follows the unkeyed sponge/hash line. This
route matches the role of the authentication tag: \TW{}'s committing statements
are consequences of tag hash security, and sponge hash security is supported by
a mature literature. The price is the indifferentiability lift of
\Cref{app:flat-sponge-lift}: one $\Lift_c(\Ntot)$ term on the real side, and in
real-vs-ideal games a second proximity term at $\Nprim$
(\Cref{cor:lift-application}).

A direct keyed-duplex treatment would have to close two gaps. First, existing
keyed-duplex bounds are hidden-key AE bounds: the game samples the key, and the
keyed initialization material is not part of the adversary's output. The
root tag claims here are adversarial-initialization statements. In the hash and
committing games, the adversary may choose the key, nonce, associated data,
message, or competing opening. Thus the proof must handle adversarially chosen
initialization transcripts, not only hidden-key encryption and verification
queries.

Second, the available keyed-duplex AE interface is not a direct substitute for a
byte-string AEAD with arbitrary final-block length. The modern bounds
synthesized by \cite{Men23} use a duplex call phased as permutation, squeeze,
then absorb or overwrite (\cite[Sec.~3.4, Table~1]{Men23}). In the
authenticated-encryption application, this interface is instantiated as
\textsf{MonkeySpongeWrap} \cite[Alg.~8]{Men23}. As written, however, the
overwrite interface does not reconstruct the missing padding bits of a truncated
final ciphertext block on decryption: encryption absorbs the padded final
plaintext block after squeezing and then truncates the resulting ciphertext to
the message length, while decryption receives only this truncated ciphertext,
pads it independently, and uses the overwrite branch to set the rate. The
decryption transcript can therefore diverge before tag recomputation.

\TW{} instead stays with the BDPV11 duplex phasing: a call absorbs the actual
input string, padded by the duplex rule, and then returns the requested output.
At the stream layer this supports arbitrary length messages by splitting them
into bounded duplex blocks, with the final short ciphertext block absorbed as it
is rather than reconstructed through an overwrite branch (\Cref{sec:tw128}). A
direct keyed-duplex theorem for \TW{} would therefore need both this BDPV11
phasing and adversarial-initialization hash and committing games. We retain the
indifferentiability lift because it covers these claims uniformly and already
meets the $128$-bit target
(\Cref{sec:bound-discussion}).

\paragraph{Beyond-birthday security and leakage.}
All bounds are birthday-type on the $256$-bit capacity, delivering the
$128$-bit target but no more; a beyond-birthday analysis is open. The model is
also single-stage and leakage-free: side-channel resistance is addressed only at
the implementation level, through the data-independent execution properties of
\Cref{sec:impl}, and is not part of the provable-security claims.
